\chapter{Esercizi: Introduzione all'informatica}
\label{chap:esercizi-introduzione-informatica}

\begin{exercise}
\label{ex:fi-intro-complessita}
Che cosa si intende per complessità di un algoritmo e perché viene
studiata?
\end{exercise}

La \emph{complessità computazionale} esprime le risorse necessarie
all'esecuzione di un algoritmo, principalmente tempo e memoria, in
funzione della dimensione dell'ingresso. Viene studiata per confrontare
algoritmi che risolvono lo stesso problema e valutarne l'efficienza
al crescere dei dati, indipendentemente dalle sole misurazioni su una
macchina particolare (\cref{sec:fi-algoritmi,chap:complessita}).

\begin{exercise}
\label{ex:fi-intro-turing-definizione}
Che cos'è la macchina di Turing e da quali elementi è composta?
\end{exercise}

La \emph{macchina di Turing} è un modello astratto di calcolo, indipendente
dalla realizzazione fisica di un calcolatore. Comprende un nastro
teoricamente illimitato suddiviso in celle, una testina di lettura e
scrittura e un controllo dotato di un numero finito di stati. Il suo
comportamento è determinato da regole di transizione
(\cref{sec:fi-macchina-turing,fig:fi-nastro-turing}).

\begin{exercise}
\label{ex:fi-intro-contributi-teorici}
Quale contributo hanno dato Turing, Church, Chomsky e Shannon allo
sviluppo teorico dell'informatica?
\end{exercise}

Turing e Church formalizzarono il concetto di calcolo effettivo e
stabilirono limiti alla decidibilità dei problemi. Chomsky introdusse
una classificazione dei linguaggi formali in base alla loro capacità
espressiva. Shannon sviluppò una teoria quantitativa dell'informazione,
introducendo misure come il contenuto informativo e l'entropia
(\cref{sec:fi-macchina-turing,sec:fi-linguaggi-esecuzione,sec:fi-informatica-calcolo}).

\begin{exercise}
\label{ex:fi-intro-algoritmo-programma}
Che cosa si intende per algoritmo e perché è necessario individuarlo prima
di scrivere un programma?
\end{exercise}

Un \emph{algoritmo} descrive un procedimento mediante istruzioni
non ambigue ed eseguibili, che termina per ogni ingresso ammissibile.
La progettazione precede l'implementazione perché occorre definire
il problema, individuare un procedimento risolutivo e verificarne la
correttezza. Un programma traduce tale procedimento in un linguaggio
di programmazione (\cref{def:fi-algoritmo,sec:fi-algoritmi}).

\begin{exercise}
\label{ex:fi-intro-problema-decisione}
Qual era il problema della decisione posto nell'ambito del formalismo
di Hilbert?
\end{exercise}

Il problema chiedeva se esistesse un procedimento generale capace
di stabilire la validità delle formule della logica del primo ordine.
Church e Turing mostrarono che un tale procedimento non esiste.
Il significato di \emph{decidibilità} è precisato nella
\cref{def:fi-decidibilita}.

\begin{exercise}
\label{ex:fi-intro-church-turing}
Quale contributo diedero Turing e Church alla teoria degli algoritmi
e al concetto di calcolabilità?
\end{exercise}

Church e Turing fornirono formalizzazioni differenti della nozione
di procedimento effettivo e dimostrarono l'esistenza di problemi
indecidibili. La \emph{tesi di Church--Turing} identifica la nozione
intuitiva di calcolo effettivo con quella di calcolo eseguibile da
una macchina di Turing (\cref{sec:fi-macchina-turing}).

\clearpage

\begin{exercise}
\label{ex:fi-intro-turing-operazioni}
Da quali elementi è composta una macchina di Turing e quali
operazioni fondamentali può eseguire?
\end{exercise}

Il nastro, la testina e il controllo a stati finiti sono descritti
nell'\cref{ex:fi-intro-turing-definizione}. A ogni passo la macchina
legge il simbolo della cella corrente e, secondo la regola di transizione,
scrive un simbolo, sposta la testina e aggiorna lo stato.
Una \emph{configurazione} comprende contenuto del nastro, posizione
della testina e stato del controllo; l'esecuzione può arrestarsi oppure
proseguire indefinitamente (\cref{sec:fi-macchina-turing}).

\begin{exercise}
\label{ex:fi-intro-importanza-turing}
Perché la macchina di Turing è considerata importante per la nascita
dell'informatica moderna?
\end{exercise}

La macchina di Turing fornisce un modello preciso di calcolo,
indipendente dall'hardware, con cui studiare quali problemi siano
risolvibili algoritmicamente. L'esistenza di una macchina universale
formalizza la possibilità di simulare programmi diversi; l'indecidibilità
del problema dell'arresto ne evidenzia i limiti teorici
(\cref{sec:fi-macchina-turing,prop:fi-problema-arresto}).

\begin{exercise}
\label{ex:fi-intro-shannon-bit}
Qual è il contributo di Shannon alla teoria dell'informazione e che
cosa rappresenta 1 bit?
\end{exercise}

Shannon introdusse una misura quantitativa dell'informazione legata
alla probabilità degli eventi. Un evento di probabilità \(p\) fornisce
\(I(p)=-\log_2 p\) bit; un esito fra due alternative equiprobabili
fornisce quindi un bit. L'\emph{entropia} misura il contenuto informativo
medio di una sorgente. Il bit come unità di informazione va distinto
dalla cifra binaria (\cref{sec:fi-informatica-calcolo}).

\begin{exercise}
\label{ex:fi-intro-significato-informatica}
Che cosa significa ``informatica'' e in che senso riguarda il
trattamento automatico dell'informazione?
\end{exercise}

L'informatica studia i metodi per rappresentare, memorizzare,
elaborare e trasmettere l'informazione mediante procedimenti
automatici. Il trattamento è \emph{automatico} quando l'elaborazione
segue istruzioni definite ed eseguibili da un calcolatore
(\cref{chap:introduzione-informatica,sec:fi-informatica-calcolo}).

\begin{exercise}
\label{ex:fi-intro-calcolatore}
Che cos'è un computer e quali funzioni fondamentali svolge?
\end{exercise}

Un \emph{calcolatore} è un dispositivo programmabile che memorizza
dati, esegue istruzioni e comunica con l'esterno. L'\emph{hardware}
comprende le componenti fisiche; il \emph{software} comprende i
programmi che ne determinano il comportamento
(\cref{sec:fi-informatica-calcolo}).

\begin{exercise}
\label{ex:fi-intro-von-neumann-bus}
Quali sono i blocchi principali del modello di von Neumann e quali
sono i vantaggi e gli svantaggi dell'uso del bus?
\end{exercise}

Il modello comprende CPU, memoria centrale e interfacce di ingresso
e uscita, collegate mediante bus per dati, indirizzi e controllo.
La memoria conserva sia istruzioni sia dati. I collegamenti condivisi
semplificano l'organizzazione del sistema, ma richiedono di coordinare
gli accessi e possono costituire un collo di bottiglia fra CPU e memoria
(\cref{sec:fi-architettura-calcolatore,fig:fi-von-neumann}).

\begin{exercise}
\label{ex:fi-intro-componenti-cpu}
Quali sono le componenti logiche principali della CPU e che ruolo
hanno ALU, unità di controllo e registri?
\end{exercise}

L'\emph{ALU} esegue operazioni aritmetiche e logiche; l'\emph{unità
di controllo} coordina l'esecuzione delle istruzioni; i \emph{registri}
conservano temporaneamente dati e informazioni di controllo. Fra
questi, il \emph{program counter} contiene l'indirizzo della prossima
istruzione (\cref{sec:fi-architettura-calcolatore}).

\begin{exercise}
\label{ex:fi-intro-clock-cpu}
Che cos'è il clock della CPU e perché la sua frequenza è una
caratteristica importante?
\end{exercise}

Il \emph{clock} scandisce le operazioni delle componenti sincrone.
Se \(T\) è il periodo, la frequenza è \(f=1/T\) e indica il numero
di cicli al secondo. È un parametro rilevante per le prestazioni,
ma non le determina da solo: un'istruzione può richiedere più cicli
e il lavoro svolto per ciclo dipende dall'architettura
(\cref{sec:fi-architettura-calcolatore}).

\begin{exercise}
\label{ex:fi-intro-pipeline}
Che cos'è l'architettura pipeline e quali sono gli stadi tipici
di una pipeline a 5 fasi?
\end{exercise}

La \emph{pipeline} suddivide l'esecuzione delle istruzioni in stadi,
permettendo di elaborare contemporaneamente fasi di istruzioni diverse.
Nel modello della \cref{fig:fi-pipeline} gli stadi sono prelievo (F),
decodifica (D), recupero degli operandi (O), esecuzione (E) e
scrittura del risultato (W). La pipeline aumenta idealmente il numero
di istruzioni completate per unità di tempo; dipendenze e conflitti
sulle risorse possono introdurre attese.

\begin{exercise}
\label{ex:fi-intro-memoria-gerarchia}
Che differenza c'è tra memoria primaria e memoria secondaria, e come
si collocano RAM, cache e memoria di massa?
\end{exercise}

La \emph{memoria primaria}, in particolare la RAM, conserva dati e
istruzioni in uso ed è normalmente volatile. La \emph{memoria
secondaria}, come SSD e HDD, offre conservazione persistente.
La \emph{cache} si colloca fra CPU e RAM nella gerarchia di memoria:
è più rapida e meno capiente della RAM e conserva copie di dati e
istruzioni (\cref{sec:fi-memoria-archiviazione,fig:fi-gerarchia-memoria}).

\begin{exercise}
\label{ex:fi-intro-gpu-parallelismo}
Perché le GPU sono adatte al calcolo parallelo e in quali ambiti
vengono usate oggi oltre alla grafica?
\end{exercise}

Le GPU dispongono di numerose unità di calcolo adatte a eseguire
operazioni simili su dati differenti in parallelo. Risultano quindi
vantaggiose quando il lavoro può essere suddiviso in molte operazioni
indipendenti, come nelle simulazioni numeriche, nell'elaborazione di
immagini e nell'addestramento di reti neurali
(\cref{sec:fi-architettura-calcolatore}).

\begin{exercise}
\label{ex:fi-intro-cuda-thread}
Che cos'è CUDA e in che modo permette di programmare una GPU usando
migliaia di thread?
\end{exercise}

\emph{CUDA} è una piattaforma NVIDIA per il calcolo su GPU. La CPU
(\emph{host}) avvia sulla GPU (\emph{device}) funzioni dette
\emph{kernel}, eseguite da thread organizzati in blocchi e griglie.
I thread di uno stesso blocco possono collaborare mediante memoria
condivisa e sincronizzazione; l'esecuzione segue il modello SIMT.
È possibile avviare migliaia di thread, che non devono necessariamente
essere tutti attivi nello stesso istante
(\cref{sec:fi-architettura-calcolatore}).
